\section{Experiments}\label{sec:experiments}

We evaluate four claims: (1) on clean benchmarks, multiple methods discriminate descendants from independents, but with different trade-offs (Section~\ref{sec:exp-baselines}); (2) the signature is robust to function-preserving reparameterization where simpler baselines fail (Section~\ref{sec:exp-laundering}); (3) the signature survives common post-training transformations (Section~\ref{sec:exp-posttraining}); and (4) behavioral similarity does not imply weight ancestry (Section~\ref{sec:exp-distill}).

\subsection{Baseline Comparison}\label{sec:exp-baselines}

\paragraph{Setup.}
We evaluate on two controlled benchmarks with known ancestry. The \textbf{MLP benchmark} (depth-16, hidden=48, input=16) trains 2 root models on synthetic regression tasks, deriving 15 descendants per root via fine-tuning, task-shift fine-tuning, weight perturbation, pruning, and quantization (3 variants each). Negatives include 8 independent models and 3 distilled students per root, yielding 52 pairs (30 positive, 22 negative). The \textbf{GPT-2 benchmark} trains 8 root language models (${\sim}$35M parameters, 6 layers) on TinyStories, deriving descendants via continued pretraining, LoRA merge, pruning, and quantization. We evaluate on test-split roots (3 roots), yielding 45 pairs (21 positive, 24 negative including 3 distilled students and 21 cross-root independents).

\paragraph{Evaluation protocol.}
Each (reference, suspect) pair receives a binary label: \emph{positive} if the suspect was derived from the reference's weights, \emph{negative} otherwise. We report AUROC (threshold-free discrimination) and Gap-$Z$ (standardized margin between positive and negative score distributions).

\paragraph{Results.}
Table~\ref{tab:baselines} compares seven methods on both benchmarks before any laundering. Weight-space methods operate directly on checkpoint weights without probe data. Activation-space methods (SVCCA, CKA) and decision-boundary methods (IPGuard) require forward passes through both models.

\begin{table*}[!t]
\centering
\small
\begin{tabular}{@{}llc cc cc@{}}
\toprule
 & & & \multicolumn{2}{c}{MLP (52 pairs)} & \multicolumn{2}{c}{GPT-2 (45 pairs)} \\
\cmidrule(lr){4-5} \cmidrule(lr){6-7}
Method & Type & Data-free & AUROC & Gap-$Z$ & AUROC & Gap-$Z$ \\
\midrule
Centered Residual Signature (ours) & Weight & \cmark & \textbf{1.00} & $+53.0$ & \textbf{1.00} & $\mathbf{+31.0}$ \\
Weight Cosine             & Weight & \cmark & \textbf{1.00} & $\mathbf{+76.3}$ & \textbf{1.00} & $+30.7$ \\
Aligned Frobenius         & Weight & \cmark & \textbf{1.00} & $+72.0$ & \textbf{1.00} & $+3.9$ \\
Singular Value Distance   & Weight & \cmark & \textbf{1.00} & $+7.9$ & 0.73 & $+0.4$ \\
SVCCA~\citep{raghu2017svcca}             & Activation & \xmark & \textbf{1.00} & $+45.4$ & 0.99 & $+3.7$ \\
CKA~\citep{kornblith2019similarity}               & Activation & \xmark & 0.83 & $+8.6$ & 0.86 & $+1.9$ \\
IPGuard~\citep{cao2021ipguard}           & Decision & \xmark & 0.70 & $-3.5$ & 0.91 & $+2.1$ \\
\bottomrule
\end{tabular}
\caption{Lineage detection baselines on MLP and GPT-2 benchmarks (no laundering). Data-free methods require no probe samples. On clean benchmarks, multiple weight-space methods achieve perfect discrimination; the proposed method's advantage emerges under reparameterization (Section~\ref{sec:exp-laundering}).}
\label{tab:baselines}
\end{table*}

On both benchmarks, weight-space methods---ours, weight cosine, and aligned Frobenius---achieve AUROC$=$1.0. Weight cosine shows the largest Gap-$Z$ margin on MLP ($+76.3$), but this advantage disappears under function-preserving reparameterization (Section~\ref{sec:exp-laundering}). Singular-value distance degrades on GPT-2 (AUROC$=$0.73), likely due to the higher-dimensional weight matrices. Activation-space methods (CKA, SVCCA) and IPGuard show consistent but imperfect performance across both benchmarks.

The key result: on unlaundered checkpoints, multiple weight-space methods achieve perfect discrimination. Simple baselines solve clean cases. The proposed method's advantage emerges under reparameterization invariance and partial-inheritance regimes.

\subsection{Robustness Under Checkpoint Laundering}\label{sec:exp-laundering}

Clean-benchmark performance measures proximity in a shared parameterization. A more stringent test asks whether the verifier captures ancestry independently of nuisance coordinate choices. We introduce \emph{function-preserving laundering}: transformations that leave the network's input--output map exactly unchanged while altering parameter coordinates.

Per-block permutation (P) permutes hidden units ($W_{\mathrm{in}} \leftarrow P W_{\mathrm{in}}$, $W_{\mathrm{out}} \leftarrow W_{\mathrm{out}} P^\top$); reciprocal rescaling (D) scales them ($W_{\mathrm{in}} \leftarrow D W_{\mathrm{in}}$, $W_{\mathrm{out}} \leftarrow W_{\mathrm{out}} D^{-1}$). These preserve the branch product $M = W_{\mathrm{out}} W_{\mathrm{in}}$ exactly ($P^\top P = I$, $D^{-1} D = I$), so our signature is invariant by construction.

\paragraph{MLP benchmark.}
We evaluate on the 52-pair MLP benchmark with six laundering variants: none, P, D-mild (log-uniform $[0.5,2]$), D-strong (log-uniform $[0.1,10]$), PD, and PDFT (P+D then 5 epochs fine-tuning). All variants pass a hard function-preservation gate ($\max|\Delta f| < 10^{-4}$).

\begin{table}[t]
\centering
\small
\setlength{\tabcolsep}{4pt}
\begin{tabular}{@{}lcccccc@{}}
\toprule
Method & none & P & D$_\mathrm{m}$ & D$_\mathrm{s}$ & PD & PDFT \\
\midrule
\textbf{Ours} & \textbf{1.0} & \textbf{1.0} & \textbf{1.0} & \textbf{1.0} & \textbf{1.0} & \textbf{1.0} \\
Re-Basin+scale & \textbf{1.0} & \textbf{1.0} & \textbf{1.0} & \textbf{1.0} & \textbf{1.0} & \textbf{1.0} \\
Raw Frobenius & \textbf{1.0} & 0.50 & \textbf{1.0} & 0.00 & 0.00 & 0.00 \\
Raw SVD dist & \textbf{1.0} & \textbf{1.0} & 0.00 & 0.00 & 0.00 & 0.00 \\
Raw weight cos & \textbf{1.0} & 0.86 & \textbf{1.0} & \textbf{1.0} & 0.80 & 0.80 \\
SVCCA & \textbf{1.0} & \textbf{1.0} & \textbf{1.0} & \textbf{1.0} & \textbf{1.0} & \textbf{1.0} \\
CKA & 0.85 & 0.85 & 0.85 & 0.85 & 0.85 & 0.99 \\
IPGuard & 0.70 & 0.70 & 0.70 & 0.70 & 0.70 & 0.83 \\
\bottomrule
\end{tabular}
\caption{AUROC under function-preserving laundering (52 pairs). P$=$permutation, D$_\mathrm{m/s}$$=$mild/strong rescaling. Our signature maintains AUROC$=$1.000 across all variants. Raw-weight baselines collapse: Frobenius fails under P or D; SVD survives P but fails D; weight cosine degrades under PD. Re-Basin+scale matches accuracy but requires $L$ Hungarian assignments per pair; ours requires zero alignment work.}
\label{tab:laundering}
\end{table}

Table~\ref{tab:laundering} shows the signature is invariant ($|\Delta\mathcal{L}| \approx 10^{-8}$) under P, D, and PD, remaining discriminative after fine-tuning (PDFT). Raw-weight baselines collapse: Frobenius drops to 0.00 under strong D; SVD survives P but fails D; weight cosine degrades to 0.80 under PD. Behavioral methods (SVCCA, CKA, IPGuard) are invariant to reparameterization but show weaker baseline discrimination.

\paragraph{GPT-2 benchmark.}
We extend to our GPT-2 benchmark, testing hidden-unit permutation: $W_{\mathrm{fc}} \leftarrow P W_{\mathrm{fc}}$, $W_{\mathrm{proj}} \leftarrow W_{\mathrm{proj}} P^\top$ per block. Four conditions: NONE, P-SUSPECT (permute suspect only), P-BOTH (different permutations on both), P+FT (permute then fine-tune 100 steps).

\begin{table}[t]
\centering
\small
\setlength{\tabcolsep}{4pt}
\begin{tabular}{@{}lcccc@{}}
\toprule
Method & NONE & P-SUS & P-BOTH & P+FT \\
\midrule
\textbf{Ours} & \textbf{1.0} & \textbf{1.0} & \textbf{1.0} & \textbf{1.0} \\
Raw weight cos & 1.0 & 1.0 & 1.0 & 1.0 \\
SVD distance & 0.76 & 0.76 & 0.76 & 0.75 \\
\bottomrule
\end{tabular}
\caption{AUROC under hidden-unit permutation on GPT-2 benchmark (92 pairs). AUROC is preserved, but score magnitudes differ: weight cosine drops 97\% (0.977$\to$0.033) under P-BOTH while ours remains at 0.970.}
\label{tab:gpt2-laundering}
\end{table}

Table~\ref{tab:gpt2-laundering} shows all methods maintain AUROC$=$1.0, but score magnitudes reveal the difference. Under P-BOTH, weight cosine scores collapse from 0.977 to 0.033---a 97\% drop---while our signature remains at 0.970. The margin compression suggests sensitivity to additional perturbations even when ranking is preserved.

\paragraph{Public LLM validation.}
We replicate across four model families (Table~\ref{tab:laundering-llm}). On each base$\to$fine-tune pair, P permutes $W_{\mathrm{gate}}$/$W_{\mathrm{up}}$ rows and $W_{\mathrm{down}}$ columns. Our score is invariant ($\Delta\mathcal{L} < 10^{-7}$); weight cosine collapses under P (0.95$\to$0.003). Unrelated controls stay at the ${\leq}10^{-4}$ null.

\begin{table}[t]
\centering
\small
\begin{tabular*}{\columnwidth}{@{\extracolsep{\fill}}llcrr@{}}
\toprule
Base & Pair & Var & $\mathcal{L}$ & W.cos \\
\midrule
\multirow{4}{*}{\shortstack[l]{LLaMA-2\\7B}}
  & chat & P & \textbf{.995} & .002 \\
  & chat & D$_{\mathrm{m}}$ & \textbf{.995} & .949 \\
  & OpenLLaMA$^\dagger$ & P & $5{\times}10^{-5}$ & .000 \\
  & OpenLLaMA$^\dagger$ & D$_{\mathrm{m}}$ & $5{\times}10^{-5}$ & .000 \\
\midrule
\multirow{4}{*}{\shortstack[l]{LLaMA-3\\8B}}
  & Inst & P & \textbf{.995} & .003 \\
  & Inst & D$_{\mathrm{m}}$ & .995 & \textbf{.954} \\
  & Mistral$^\dagger$ & P & $0$ & .000 \\
  & Mistral$^\dagger$ & D$_{\mathrm{m}}$ & $0$ & .000 \\
\midrule
\multirow{2}{*}{\shortstack[l]{Mistral\\7B}}
  & Inst & P & \textbf{.952} & .003 \\
  & Inst & D$_{\mathrm{m}}$ & .952 & \textbf{.937} \\
\midrule
\multirow{2}{*}{\shortstack[l]{Qwen2.5\\7B}}
  & Inst & P & \textbf{.999} & .003 \\
  & Inst & D$_{\mathrm{m}}$ & .999 & \textbf{.951} \\
\bottomrule
\end{tabular*}
\caption{Function-preserving laundering on public checkpoints. $\mathcal{L}$ is invariant ($\Delta < 10^{-7}$); weight cosine collapses under P. $^\dagger$Unrelated controls.}
\label{tab:laundering-llm}
\end{table}

\paragraph{Layer grafts and partial inheritance.}
On layer grafts and linear merges (Appendix~\ref{app:harder-bench}), $\mathcal{L}$ tracks the fraction of inherited blocks with Spearman $\rho \approx 0.96$, measuring partial overlap rather than merely crossing a threshold. Singular-value distance fails the merge regime (AUROC$=$0.448), and IPGuard inverts on grafts (AUROC$=$0.05), ranking higher-overlap suspects as \emph{less} similar.

\paragraph{Public checkpoint case study.}
On unlaundered public checkpoints (Table~\ref{tab:real-llm}), we compare LLaMA-2-7B-base against 3 documented derivatives and 1 independently trained model (OpenLLaMA-7B). The three derivatives produce high scores ($\mathcal{L}$: 0.995, 0.996, 0.336); the independent model scores $4 \times 10^{-4}$, a ${\sim}1000\times$ gap.

\begin{table}[t]
\centering\footnotesize
\setlength{\tabcolsep}{3pt}
\renewcommand{\arraystretch}{1.1}
\begin{tabular}{@{}llrrl@{}}
\toprule
Pair (base vs.) & Kind & $\mathcal{L}$ & W.\,cos & Verdict \\
\midrule
Llama-2-7B-chat   & RLHF          & $\mathbf{0.995}$ & $0.995$ & REL. \\
Vicuna-7B         & instruct      & $\mathbf{0.996}$ & $0.996$ & REL. \\
CodeLlama-7B      & code PT       & $0.336$ & $\mathbf{0.414}$ & REL. \\
OpenLLaMA-7B      & scratch$^\dagger$ & $\mathbf{4\mathrm{e}{-}4}$ & $0.089$ & UNR. \\
\bottomrule
\end{tabular}
\caption{Public checkpoint case study (LLaMA-2 family). Scores are consistent with documented relationships. $^\dagger$OpenLLaMA shares architecture but is trained from scratch.}
\label{tab:real-llm}
\end{table}

CodeLlama's lower score ($\mathcal{L}=0.336$) reflects extensive continued pretraining on code data. This case study provides real-world plausibility but has limited calibration power: only one confirmed independent model exists for this architecture.

\subsection{Robustness to Post-Training Techniques}\label{sec:exp-posttraining}

We test whether the signature survives common post-training transformations that modify weights while preserving (or approximately preserving) model behavior. We train 8 GPT-2-style transformers from scratch on TinyStories~\citep{eldan2023tinystories} with known ancestry. Each root (${\sim}$35M parameters: 6 layers, $d{=}384$, 50K samples, 3 epochs) is allocated to calibration ($n{=}2$), development ($n{=}3$), or test ($n{=}3$) splits. Per root we generate 7 descendants: 1 fine-tuned (1 epoch continued pretraining), 1 LoRA-merged (rank-8), 3 pruned (30/50/70\% sparsity), and 2 quantized (INT8/INT6).

\begin{table}[t]
\centering\small
\setlength{\tabcolsep}{4pt}
\begin{tabular}{@{}lccc@{}}
\toprule
Transformation & $n$ & Mean $\mathcal{L}$ & Min $\mathcal{L}$ \\
\midrule
Quantized (INT8/6) & 6 & 0.999 & 0.999 \\
LoRA merge (rank-8) & 3 & 0.998 & 0.996 \\
Fine-tuned (1 epoch) & 3 & 0.980 & 0.980 \\
Pruned (30--70\%) & 9 & 0.937 & 0.855 \\
\midrule
Distilled student & 3 & 0.002 & 0.001 \\
Independent & 21 & 0.003 & 0.000 \\
\bottomrule
\end{tabular}
\caption{Lineage scores under post-training transformations (GPT-2 benchmark, 3 test roots, 45 pairs). AUROC$=$1.000 with perfect separation. The minimum descendant score (0.855 at 70\% sparsity) exceeds the maximum non-descendant score (0.004) by $200\times$.}
\label{tab:gpt2-survival}
\end{table}

Table~\ref{tab:gpt2-survival} reports lineage scores on the held-out test split. Quantization and LoRA preserve the signal almost perfectly ($\mathcal{L}{\geq}0.996$); fine-tuning and pruning degrade it but remain well above the null ($\mathcal{L}{\geq}0.855$). Both distilled students and independently trained models score near zero, correctly classified as unrelated. The $200\times$ gap between minimum descendant and maximum non-descendant scores suggests robustness to these common transformations.

\subsection{Behavioral Similarity Does Not Imply Weight Ancestry}\label{sec:exp-distill}

A key question is whether behavioral similarity---output agreement between models---implies weight-level ancestry. We test this using the distillation control from our GPT-2 benchmark.

\begin{table}[t]
\centering\small
\setlength{\tabcolsep}{4pt}
\begin{tabular}{@{}lccc@{}}
\toprule
Comparison & Top-1 (\%) & KL & $\mathcal{L}$ \\
\midrule
Teacher--Distilled & 79.0 & 0.16 & 0.002 \\
Teacher--Independent & 77.6 & 0.25 & 0.003 \\
\bottomrule
\end{tabular}
\caption{Behavioral similarity vs.\ lineage score. Models trained on the same architecture and corpus already achieve 77.6\% top-1 agreement. Distillation increases agreement modestly (+1.4 points) and reduces KL divergence ($0.25 \to 0.16$), but neither independently trained nor distilled models exhibit weight ancestry---both score $\mathcal{L} \approx 0.002$.}
\label{tab:distill-quality}
\end{table}

Table~\ref{tab:distill-quality} shows that models trained on the same architecture and corpus already exhibit substantial behavioral agreement: independent models achieve 77.6\% top-1 agreement. Distillation increases this modestly to 79.0\% (+1.4 points) and reduces KL divergence from 0.25 to 0.16. Yet both independently trained models and distilled students score near zero in lineage ($\mathcal{L} \approx 0.002$), correctly classified as unrelated.

\textbf{Interpretation.} This experiment illustrates the conceptual distinction between behavioral similarity and weight-level lineage. A model trained from scratch is unrelated under our task definition, regardless of how closely it imitates another model's outputs. The result validates the task definition rather than differentiating our method from baselines: all evaluated lineage methods correctly reject these distilled students, because the current distillation procedure produces only modest behavioral agreement (+1.4 points over independent training).

\textbf{Limitation.} Stronger distillation that achieves near-perfect imitation would be needed to test whether behavioral or activation-space baselines produce false lineage evidence under closer imitation. The current experiment establishes that output similarity is not equivalent to weight inheritance, but does not establish that our method outperforms behavioral baselines on challenging distillation scenarios.
